% --- CORRECCION: Forzar interpretacion raw de la entrada ---
\UseRawInputEncoding
\documentclass[aspectratio=169, 10pt, xcolor=table]{beamer}

% --- Configuracion del Tema ---
\usetheme{Madrid}
\usecolortheme{dolphin}

% --- Paquetes necesarios ---
\usepackage[utf8]{inputenc}
\usepackage[spanish]{babel}
\usepackage{amsmath, amssymb}
\usepackage{booktabs}
\usepackage{graphicx}
\usepackage{listings}
\usepackage{tikz}
\usetikzlibrary{arrows.meta, positioning, shapes.geometric}
\usepackage{etoolbox}

% --- Ruta de Imagenes ---
\graphicspath{{./imagenes/}}

% --- Estilo para codigo Python ---
\definecolor{codegreen}{rgb}{0,0.6,0}
\definecolor{codegray}{rgb}{0.5,0.5,0.5}
\definecolor{codepurple}{rgb}{0.58,0,0.82}
\definecolor{backcolour}{rgb}{0.95,0.95,0.92}
\lstset{
    language=Python,
    backgroundcolor=\color{backcolour},
    commentstyle=\color{codegreen},
    keywordstyle=\color{blue},
    numberstyle=\tiny\color{codegray},
    stringstyle=\color{codepurple},
    basicstyle=\ttfamily\scriptsize,
    breaklines=true,
    frame=single,
    showstringspaces=false,
    literate={á}{{\'a}}1 {é}{{\'e}}1 {í}{{\'i}}1 {ó}{{\'o}}1 {ú}{{\'u}}1 {ñ}{{\~n}}1 {ü}{{\"u}}1
}

% --- Pie de pagina personalizado ---
\makeatletter
\setbeamertemplate{footline}{
  \leavevmode%
  \hbox{%
  \begin{beamercolorbox}[wd=.333333\paperwidth,ht=2.25ex,dp=1ex,center]{author in head/foot}%
    \usebeamerfont{author in head/foot}\insertshortauthor
  \end{beamercolorbox}%
  \begin{beamercolorbox}[wd=.333333\paperwidth,ht=2.25ex,dp=1ex,center]{title in head/foot}%
    \usebeamerfont{title in head/foot}Sesion 9
  \end{beamercolorbox}%
  \begin{beamercolorbox}[wd=.333333\paperwidth,ht=2.25ex,dp=1ex,right]{date in head/foot}%
    \usebeamerfont{date in head/foot}CALCULO Y ALGEBRA LINEAL \hfill \insertframenumber/\inserttotalframenumber \hspace*{0.3cm}
  \end{beamercolorbox}}%
  \vskip0pt%
}
\makeatother

% --- Datos de la portada ---
\title[SESION 9 CALCULO]{\textbf{Sesion 9}}
\subtitle{Automatizando el Gradiente: PyTorch/TensorFlow y Descenso de Gradiente}
\author[Prof. C. Sanchez]{Carlos Cesar Sanchez Coronel}
\institute{\textbf{CALCULO Y ALGEBRA LINEAL PARA IA}}
\date{}

% --- Eliminar barra de navegacion ---
\setbeamertemplate{navigation symbols}{}

% --- Indice al inicio de cada seccion ---
\AtBeginSection[]{
  \begin{frame}
    \frametitle{Contenido}
    \tableofcontents[currentsection]
  \end{frame}
}

\begin{document}

% --- Portada ---
\begin{frame}
    \titlepage
\end{frame}

% --- Indice general ---
\begin{frame}{Contenido}
    \tableofcontents
\end{frame}

% ============================================================================
% SECCION 1: DESCENSO DE GRADIENTE: RECORDATORIO Y VARIANTES
% ============================================================================
\section{Descenso de Gradiente: Recordatorio y Variantes}

\begin{frame}{El algoritmo basico}
    \begin{block}{Descenso de gradiente}
        \[
        \mathbf{w}_{t+1} = \mathbf{w}_t - \eta \nabla \mathcal{L}(\mathbf{w}_t)
        \]
        donde $\eta$ es la tasa de aprendizaje y $\nabla \mathcal{L}(\mathbf{w}_t)$ es el gradiente de la perdida.
    \end{block}
\end{frame}

\begin{frame}{Epocas, batches y variantes}
    \begin{columns}[t]
        \column{0.33\textwidth}
        \textbf{Batch GD}
        \begin{itemize}
            \item Usa todo el dataset.
            \item Estable, costoso.
        \end{itemize}
        \column{0.33\textwidth}
        \textbf{SGD (estocastico)}
        \begin{itemize}
            \item Una muestra por iteracion.
            \item Rapido, ruidoso.
        \end{itemize}
        \column{0.33\textwidth}
        \textbf{Mini-batch GD}
        \begin{itemize}
            \item Subconjunto (batch).
            \item Compromiso ideal.
        \end{itemize}
    \end{columns}
    \vspace{0.3cm}
    Una \textbf{epoca} = recorrer todo el dataset una vez.
\end{frame}

% ============================================================================
% SECCION 2: DIFERENCIACION AUTOMATICA Y GRAFOS COMPUTACIONALES
% ============================================================================
\section{Diferenciacion Automatica y Grafos Computacionales}

\begin{frame}{Comparativa de metodos de diferenciacion}
    \begin{table}
        \centering
        \begin{tabular}{l|c|c|c}
            \toprule
            Metodo & Exactitud & Costo & Ventaja \\
            \midrule
            Numerico (dif. finitas) & Baja & Medio & Facil implementacion \\
            Simbolico (SymPy) & Exacta & Muy alto & Expresion analitica \\
            Automatico (Autodiff) & Exacta & Bajo & Eficiencia en $n$ vars \\
            \bottomrule
        \end{tabular}
    \end{table}
\end{frame}

\begin{frame}{Grafos computacionales}
    \begin{itemize}
        \item Representacion de operaciones como grafo dirigido aciclico (DAG).
        \item Cada nodo: operacion o variable.
        \item Permite aplicar regla de la cadena automaticamente.
    \end{itemize}
    \centering
    \begin{tikzpicture}[node distance=1.2cm, auto]
        \node (x) {$x$};
        \node (w) [right of=x] {$w$};
        \node (b) [right of=w] {$b$};
        \node (mul) [below of=w] {$\times$};
        \node (sum) [below of=b] {$+$};
        \node (y) [below of=sum] {$y$};
        \draw[->] (x) -- (mul);
        \draw[->] (w) -- (mul);
        \draw[->] (mul) -- (sum);
        \draw[->] (b) -- (sum);
        \draw[->] (sum) -- (y);
    \end{tikzpicture}
\end{frame}

\begin{frame}{Anatomia de una funcion en grafos: $f(x)=\ln(x^2+1)$}
    \begin{columns}[t]
        \column{0.5\textwidth}
        \textbf{Forward pass (primals)}
        \begin{align*}
            v_1 &= x^2 = 4 \\
            v_2 &= v_1 + 1 = 5 \\
            v_3 &= \ln(v_2) \approx 1.6094
        \end{align*}
        \column{0.5\textwidth}
        \textbf{Backward pass (sensibilidad)}
        \begin{align*}
            \frac{\partial f}{\partial v_3} &= 1 \\
            \frac{\partial f}{\partial v_2} &= 1\cdot \frac{1}{v_2}=0.2 \\
            \frac{\partial f}{\partial v_1} &= 0.2\cdot 1 = 0.2 \\
            \frac{df}{dx} &= 0.2\cdot 2x = 0.8
        \end{align*}
    \end{columns}
\end{frame}

\begin{frame}{Reglas del backward pass}
    \begin{enumerate}
        \item Empezar con 1 al final del grafo.
        \item En cada flecha hacia atras, multiplicar el gradiente acumulado por la derivada local de la operacion.
        \item Usar valores del forward pass para evaluar derivadas locales.
    \end{enumerate}
\end{frame}

% ============================================================================
% SECCION 3: IMPLEMENTACION MANUAL CON NUMPY
% ============================================================================
\section{Implementacion Manual con NumPy}

\begin{frame}[fragile]{Generacion de datos sinteticos}
    \begin{lstlisting}
import numpy as np
import matplotlib.pyplot as plt

np.random.seed(42)
X = 2 * np.random.rand(100, 1)
y = 2 + 3 * X + np.random.randn(100, 1)
    \end{lstlisting}
\end{frame}

\begin{frame}[fragile]{Clase LinearRegressionSGD}
    \begin{lstlisting}
class LinearRegressionSGD:
    def __init__(self, lr=0.01, epochs=100, batch_size=32):
        self.lr = lr
        self.epochs = epochs
        self.batch_size = batch_size
        self.w = None
        self.b = None
        self.loss_history = []

    def fit(self, X, y):
        n_samples, n_features = X.shape
        self.w = np.random.randn(n_features, 1) * 0.01
        self.b = 0.0
        for epoch in range(self.epochs):
            indices = np.random.permutation(n_samples)
            X_shuffled = X[indices]
            y_shuffled = y[indices]
            epoch_loss = 0.0
            for i in range(0, n_samples, self.batch_size):
                X_batch = X_shuffled[i:i+self.batch_size]
                y_batch = y_shuffled[i:i+self.batch_size]
                y_pred = X_batch @ self.w + self.b
                loss = np.mean((y_pred - y_batch)**2)
                grad_w = 2 * X_batch.T @ (y_pred - y_batch) / len(X_batch)
                grad_b = 2 * np.mean(y_pred - y_batch)
                self.w -= self.lr * grad_w.reshape(-1,1)
                self.b -= self.lr * grad_b
                epoch_loss += loss * len(X_batch)
            self.loss_history.append(epoch_loss / n_samples)
    \end{lstlisting}
\end{frame}

% ============================================================================
% SECCION 4: INTRODUCCION A PYTORCH
% ============================================================================
\section{Introduccion a PyTorch}

\begin{frame}[fragile]{Tensores en PyTorch}
    \begin{lstlisting}
import torch

x = torch.tensor([1.0, 2.0, 3.0])
w = torch.tensor([0.5, -0.2, 0.1], requires_grad=True)
b = torch.tensor(0.0, requires_grad=True)
y = (x * w).sum() + b
print(y.item())
    \end{lstlisting}
\end{frame}

\begin{frame}[fragile]{Calculo de gradientes con autograd}
    \begin{lstlisting}
x = torch.tensor(2.0, requires_grad=True)
y = torch.log(x**2 + 1)
y.backward()
print(f"Gradiente en x=2: {x.grad}")  # 0.8
    \end{lstlisting}
\end{frame}

\begin{frame}[fragile]{Entrenamiento de regresion lineal con PyTorch}
    \begin{lstlisting}
import torch.nn as nn
import torch.optim as optim

X_t = torch.tensor(X, dtype=torch.float32)
y_t = torch.tensor(y, dtype=torch.float32)

model = nn.Linear(1, 1)
criterion = nn.MSELoss()
optimizer = optim.SGD(model.parameters(), lr=0.1)

dataset = torch.utils.data.TensorDataset(X_t, y_t)
dataloader = torch.utils.data.DataLoader(dataset, batch_size=16, shuffle=True)

for epoch in range(50):
    for X_batch, y_batch in dataloader:
        y_pred = model(X_batch)
        loss = criterion(y_pred, y_batch)
        optimizer.zero_grad()
        loss.backward()
        optimizer.step()
    \end{lstlisting}
\end{frame}

% ============================================================================
% SECCION 5: INTRODUCCION A TENSORFLOW / KERAS
% ============================================================================
\section{Introduccion a TensorFlow / Keras}

\begin{frame}[fragile]{Modelo secuencial en Keras}
    \begin{lstlisting}
import tensorflow as tf
from tensorflow import keras

model = keras.Sequential([
    keras.layers.Dense(1, input_shape=(1,))
])
model.compile(optimizer='sgd', loss='mse')
history = model.fit(X, y, epochs=50, batch_size=16, verbose=0)
print("Peso:", model.layers[0].weights[0].numpy()[0,0])
print("Intercepto:", model.layers[0].weights[1].numpy()[0])
    \end{lstlisting}
\end{frame}

% ============================================================================
% SECCION 6: OPTIMIZADORES MODERNOS
% ============================================================================
\section{Optimizadores Modernos}

\begin{frame}{Momentum}
    Acumula una media exponencial de gradientes:
    \[
    v_{t+1} = \beta v_t + \nabla\mathcal{L}(w_t),\quad w_{t+1}=w_t-\eta v_{t+1}
    \]
\end{frame}

\begin{frame}{Adam (Adaptive Moment Estimation)}
    Combina Momentum y RMSprop:
    \begin{align*}
        m_t &= \beta_1 m_{t-1} + (1-\beta_1)g_t \\
        v_t &= \beta_2 v_{t-1} + (1-\beta_2)g_t^2 \\
        \hat{m}_t &= \frac{m_t}{1-\beta_1^t},\quad \hat{v}_t = \frac{v_t}{1-\beta_2^t} \\
        \theta_{t+1} &= \theta_t - \eta \frac{\hat{m}_t}{\sqrt{\hat{v}_t}+\epsilon}
    \end{align*}
    $\beta_1=0.9$, $\beta_2=0.999$, $\epsilon=10^{-8}$.
\end{frame}

\begin{frame}[fragile]{Uso de Adam en PyTorch y Keras}
    PyTorch:
    \begin{lstlisting}
optimizer = optim.Adam(model.parameters(), lr=0.001)
    \end{lstlisting}
    Keras:
    \begin{lstlisting}
optimizer = keras.optimizers.Adam(learning_rate=0.001)
    \end{lstlisting}
\end{frame}

% ============================================================================
% SECCION 7: COMPARACION DE IMPLEMENTACIONES
% ============================================================================
\section{Comparacion de Implementaciones}

\begin{frame}{Tabla comparativa}
    \begin{table}
        \centering
        \begin{tabular}{l|c|c|c}
            \toprule
            Caracteristica & NumPy manual & PyTorch & TensorFlow/Keras \\
            \midrule
            Calculo de gradientes & Manual & Automatico & Automatico \\
            Soporte GPU & No & Si & Si \\
            Facilidad de prototipado & Baja & Alta & Muy alta \\
            Control detallado & Total & Total & Medio \\
            Curva de aprendizaje & Alta & Media & Baja \\
            \bottomrule
        \end{tabular}
    \end{table}
\end{frame}

% ============================================================================
% SECCION 8: NOTACION EN PAPERS CIENTIFICOS
% ============================================================================
\section{Notacion en Papers Cientificos}

\begin{frame}{Notacion comun en deep learning}
    \begin{itemize}
        \item $\theta_{t+1} = \theta_t - \eta \nabla_\theta \mathcal{L}_B(\theta_t)$: actualizacion con mini-batch $B$.
        \item $\delta^{(l)}$: error propagado hacia atras en capa $l$.
        \item $\frac{\partial \mathcal{L}}{\partial \mathbf{W}^{(l)}} = \delta^{(l)} (\mathbf{a}^{(l-1)})^\top$.
        \item Adam y otros optimizadores suelen describirse con sus ecuaciones.
    \end{itemize}
\end{frame}

% ============================================================================
% SECCION 9: EJERCICIOS PROPUESTOS
% ============================================================================
\section{Ejercicios Propuestos}

\begin{frame}{Ejercicios}
    \begin{enumerate}
        \item Implementa SGD puro (batch=1) con NumPy y observa la varianza.
        \item Modifica el codigo PyTorch para usar Adam y compara convergencia con SGD.
        \item Construye el grafo computacional de $f(x)=e^{\sin(x^2)}$ y calcula su derivada en $x=1$ paso a paso; verifica con PyTorch.
        \item Investiga el problema del \textit{vanishing gradient} y como los optimizadores modernos ayudan a mitigarlo.
        \item Lee la seccion de optimizacion del paper original de Adam e identifica la notacion.
    \end{enumerate}
\end{frame}

% ============================================================================
% SECCION 10: RESUMEN
% ============================================================================
\section{Resumen}

\begin{frame}{Lo que hemos aprendido}
    \begin{exampleblock}{Conceptos clave}
        \begin{itemize}
            \item Variantes de descenso de gradiente (batch, SGD, mini-batch).
            \item Diferenciacion automatica y grafos computacionales.
            \item Implementacion manual con NumPy.
            \item Fundamentos de PyTorch (tensores, autograd, entrenamiento).
            \item Fundamentos de TensorFlow/Keras.
            \item Optimizadores modernos: Momentum, Adam.
            \item Comparacion de frameworks.
            \item Notacion en papers.
        \end{itemize}
    \end{exampleblock}
\end{frame}

\begin{frame}{Conexion con futuros cursos}
    \begin{itemize}
        \item \textbf{Deep Learning}: redes profundas, backpropagation.
        \item \textbf{Optimizacion avanzada}: AdamW, RMSprop, etc.
        \item \textbf{Implementacion de modelos}: uso de frameworks en proyectos reales.
    \end{itemize}
\end{frame}

% --- Diapositiva final ---
\begin{frame}
    \centering
    \Huge \textbf{Gracias!}

\end{frame}

\end{document}